\section{The integral extension of the two-fibre class}

The coefficient $\Delta_{b_0}/g$ determines the coarse transgression.
The total integral group can still be a nonsplit extension.  For example,
$r=s=2$, $b_0=0$, and $m=n=-1$ give a cyclic group of order four,
although both successive quotients have order two.

Retain the group $\mathcal S_{r,s}$ in (4.9), for arbitrary $r,s>0$ and
$m,n,b_0\in\mathbb Z$.  Write
\[
 g=\gcd(r,s),\qquad R=r/g,\qquad S=s/g,\qquad
 \ell=gRS,\qquad k=\Delta_{b_0}/g.
\]
Choose integers $a,b$ with $aS+bR=1$, and put
\begin{equation}\label{eq:seifert-split}
 U=aX_0+bX_1,\qquad T=RX_0-SX_1,\qquad t=bm-an.
\end{equation}
The letters $a,b$ in this section denote B\'ezout coefficients.

\begin{theorem}[Integral extension and splitting]\label{thm:extension}
One has
\begin{equation}\label{eq:class-coordinates}
 \mathcal S_{r,s}=\mathbb ZU\oplus(\mathbb Z/g\mathbb Z)T,
 \qquad H=\ell U,\qquad \varepsilon=kU+tT.
\end{equation}
Consequently, for $E=\mathcal S_{r,s}/\mathbb Z\varepsilon$,
\begin{equation}\label{eq:extension-presentation}
 E=\langle U,T\mid gT=0,\ kU+tT=0\rangle.
\end{equation}
If $k\ne0$, this gives an exact sequence
\begin{equation}\label{eq:extension-sequence}
 0\longrightarrow\mathbb Z/g\mathbb Z
 \mathrel{\mathop{\longrightarrow}^{1\mapsto T}}E
 \mathrel{\mathop{\longrightarrow}^{U\mapsto1,\ T\mapsto0}}
 \mathbb Z/|k|\mathbb Z\longrightarrow0.
\end{equation}
It splits if and only if $\gcd(g,k)$ divides $t$.
Moreover, with $d=\gcd(r,s,m,n)$,
\[
 \gcd(g,k,t)=d,\qquad
 E\cong
 \begin{cases}
 \mathbb Z/d\mathbb Z\oplus\mathbb Z/(|\Delta_{b_0}|/d)\mathbb Z,
       &k\ne0,\\
 \mathbb Z\oplus\mathbb Z/d\mathbb Z,&k=0.
 \end{cases}
\]
Thus $E$ and $G_{b_0}$ have the same invariant factors.
This assertion alone specifies no canonical isomorphism between their generators.
\end{theorem}

\begin{proof}
Eliminate $H$ from (4.9).  The remaining relation is
$g(RX_0-SX_1)=0$.  The change of generators in
\eqref{eq:seifert-split} has determinant $-aS-bR=-1$, with inverse
\[
             X_0=SU+bT,\qquad X_1=RU-aT.
\]
It follows that $T$ has exact order $g$ and $H=rX_0=\ell U$.
Substitution in (4.10) gives \eqref{eq:class-coordinates} and
\eqref{eq:extension-presentation}.

If $k\ne0$, no nonzero multiple of the relation $kU+tT$ lies in
$\mathbb ZT$.  Thus the indicated cyclic subgroup injects, and the
quotient is $\mathbb Z/|k|\mathbb Z$.  A lift of its generator has
form $U+jT$ and is killed by $k$ precisely when
$kj\equiv t\pmod g$.  This congruence is soluble exactly under the
stated divisibility condition.

The matrix sending $(m,n)$ to $(Sm+Rn,bm-an)$ is unimodular.
Since $g\mid\ell$, this gives
\[
 \gcd(g,k,t)=\gcd(g,Sm+Rn,bm-an)=\gcd(g,m,n)=d.
\]
The relation matrix in \eqref{eq:extension-presentation} has nonzero
entry $g$, determinant of absolute value $g|k|=|\Delta_{b_0}|$,
and first determinantal divisor $d$.  Its Smith form proves both cases.
\end{proof}

Replacing $(a,b)$ by $(a+Rj,b-Sj)$ changes $U$ to $U+jT$ and changes
$t$ to $t-jk$ modulo $g$.  The presentation and the splitting condition
are therefore independent of the B\'ezout choice.
The section gauges (3.7) change $t$ by a multiple of $g$ and leave $k$ fixed.

\begin{proposition}[The coarse Leray extension]\label{prop:leray-extension}
For the effective rank-one Seifert bundle (4.3), let $M$ be its oriented
unit-circle subbundle.  If $k\ne0$, the integral Leray filtration gives
\begin{equation}\label{eq:leray-extension}
 0\longrightarrow\mathbb Z/|k|\mathbb Z
 \longrightarrow H^2(M,\mathbb Z)
 \longrightarrow\mathbb Z/g\mathbb Z\longrightarrow0.
\end{equation}
Its middle group is cyclic of order $|\Delta_{b_0}|$.
The extension splits if and only if $\gcd(g,k)=1$.
For $k=0$, the middle group is $\mathbb Z$ and the sequence is
$0\to\mathbb Z\xrightarrow{\,\pm g\,}\mathbb Z\to\mathbb Z/g\to0$.
\end{proposition}

\begin{proof}
A neighbourhood of each fibre retracts onto a circle, so
$R^jq_*\mathbb Z=0$ for $j>1$.  The base has cohomological dimension
two.  Hence the only terms contributing to total degree two are
$E_\infty^{2,0}=\operatorname{coker}d_2^{0,1}$ and
$E_\infty^{1,1}=\mathbb Z/g\mathbb Z$.  Equation (4.1) gives the
left group in \eqref{eq:leray-extension}.  Radial contraction from $Z$
to $M$ preserves the base and its Leray filtration.

Effectiveness implies $d=1$.  Theorem~3.1 and abelianity give
$H_1(M,\mathbb Z)=G_{b_0}$, and integral Poincar\'e duality gives
$H^2(M,\mathbb Z)\cong H_1(M,\mathbb Z)$.
If $k\ne0$, a direct sum of cyclic groups of orders $g$ and $|k|$ is
cyclic exactly when those orders are coprime.  Conversely, a cyclic group
of order $g|k|$ has complementary subgroups of these orders when they
are coprime.  This proves the splitting assertion.
If $k=0$, Theorem~3.1 gives the middle group $\mathbb Z$.
Its subgroup with quotient $\mathbb Z/g\mathbb Z$ is $g\mathbb Z$.
\end{proof}

The algebraic sequence \eqref{eq:extension-sequence} places the two
cyclic groups in the opposite order from \eqref{eq:leray-extension}.
Their middle groups have the same invariant factors.  No identification
of the two filtrations is asserted.
For $r=s=2$, $m=n=-1$, choose $a=1,b=0$.  Then $t=1$ and
$k=-2(b_0+1)$.  At $b_0=0$, both extensions have middle group
$\mathbb Z/4\mathbb Z$ and do not split.
At $b_0=-1$, the nonzero torsion class $\varepsilon=T$ has zero
coarse coefficient, and quotienting by it leaves $E\cong\mathbb Z$.

\section{The linking form of an effective filling}

A group presentation determines the orders of cycles.  Their linking
also uses the oriented filling slopes.  We retain this geometric data
explicitly, without selecting a convention for naming lens spaces.

Let $T^2$ have oriented basis $(x,c)$ with intersection $x\cdot c=1$.
Set $\nu=n-sb_0$ and regard
\begin{equation}\label{eq:filling-slopes}
       \mu_0=rx-mc,\qquad \mu_1=sx+\nu c
\end{equation}
as curves on $T^2$.  Suppose
$\gcd(r,m)=\gcd(s,n)=1$.  Both curves are then primitive.
Take an oriented solid torus $V_0$ whose meridian is $\mu_0$, identifying
$T^2$ with its boundary with the boundary orientation.
Glue a second solid torus $V_1$ with meridian $\mu_1$ by an
orientation-reversing boundary map.  Denote the resulting oriented
closed manifold by $M_A$.
The slope $c$ specifies its regular Seifert fibre.  The intersection
numbers with the two meridians have absolute values $r,s$.
The section relation $x_1=-x+b_0c$ identifies the second local
relation with $-\mu_1$.  Thus these fillings realize the data of (3.1),
with the stated orientation on $T^2$.
The gluing description is the two-fibre case of the solid-torus
construction in \cite[Section~2.1, pp.~25--27]{Hatcher}.

For disjoint torsion cycles $z,w$ in an oriented closed three-manifold,
our convention is
\[
 \lambda_M([z],[w])=\frac{C\cdot w}{N}\pmod{\mathbb Z},
             \qquad \partial C=Nz,\quad N\ne0.
\]
This is the symmetric torsion linking pairing
\cite[Section~1, pp.~2--3]{Hillman}.

\begin{theorem}[Two-fibre linking coefficient]\label{thm:linking}
Choose integers $\alpha,\beta$ with $m\alpha+r\beta=1$, and put
\begin{equation}\label{eq:linking-coordinates}
 \lambda_0=\alpha x+\beta c,\qquad
 q=s\beta-\nu\alpha,\qquad \Delta=rn+sm-rsb_0.
\end{equation}
Then
\[
       \mu_1=q\mu_0+\Delta\lambda_0,\qquad
       \gcd(q,\Delta)=1.
\]
If $\Delta\ne0$, the class $\gamma=[\lambda_0]$ generates
$H_1(M_A,\mathbb Z)\cong\mathbb Z/|\Delta|\mathbb Z$, and
\begin{equation}\label{eq:linking-formula}
 \lambda_{M_A}(i\gamma,j\gamma)=-\frac{qij}{\Delta}
            \pmod{\mathbb Z},\qquad
 [x]=m\gamma,\qquad[c]=r\gamma.
\end{equation}
The pairing is nonsingular.  For $\Delta=0$, one has
$H_1(M_A,\mathbb Z)\cong\mathbb Z$ and the torsion pairing is trivial.
\end{theorem}

\begin{proof}
The columns $(\mu_0,\lambda_0)$ form a unimodular matrix of determinant
$r\beta+m\alpha=1$.  Taking determinants in this basis gives
\[
 \mu_1=(s\beta-\nu\alpha)\mu_0+(r\nu+sm)\lambda_0.
\]
Primitivity of $\mu_1$ proves $\gcd(q,\Delta)=1$.
The inverse change of basis is
\[
               x=\beta\mu_0+m\lambda_0,\qquad
               c=-\alpha\mu_0+r\lambda_0.
\]
The two meridional disks kill $\mu_0$ and $\mu_1$ in first homology,
so the resulting presentation has the single relation $\Delta\gamma=0$.

Let $D_0,D_1$ be the oriented meridional disks with respective
boundaries $\mu_0,\mu_1$.  On $T^2$, the cycle
$\mu_1-q\mu_0-\Delta\lambda_0$ bounds an integral singular
2-chain $B$, because its homology class is zero.  Consequently
\[
                   C=D_1-qD_0-B,\qquad
                   \partial C=\Delta\lambda_0.
\]
Push a representative of $\lambda_0$ into the interior of $V_0$.
It is disjoint from $D_1$ and $B$ and intersects $D_0$ with sign $+1$.
The sign follows from the boundary orientation and
$\mu_0\cdot\lambda_0=1$.  Therefore $C\cdot\lambda_0=-q$,
which proves \eqref{eq:linking-formula}.  Bilinearity gives all values.
Since $q$ is a unit modulo $\Delta$, this pairing identifies the cyclic
group with its character group.  If $\Delta=0$, its first homology
is infinite cyclic and has no torsion.
\end{proof}

Replacing $(\alpha,\beta)$ by $(\alpha+rj,\beta-mj)$ replaces
$\lambda_0$ by $\lambda_0+j\mu_0$ and $q$ by $q-j\Delta$.
Thus neither $\gamma$ nor the linking value changes.
Reversing the orientation of $M_A$ negates the pairing.
A different cyclic generator multiplies its displayed coefficient by a
unit square modulo $|\Delta|$.
These rules state the full dependence on choices, including even orders.

For $(r,s;m,n;b_0)=(2,2;-1,-1;0)$, take
$(\alpha,\beta)=(-1,0)$.  Then $\Delta=-4$, $q=-1$, and
$\lambda_{M_A}(\gamma,\gamma)=-1/4$.
For $(r,s;m,n;b_0)=(2,3;1,1;0)$, take $(\alpha,\beta)=(1,0)$.
Now $\Delta=5$, $q=-1$, and the value is $1/5$.
The nonsingular symmetric form on $\mathbb Z/5\mathbb Z$ cannot
be an alternating form: every alternating bilinear form on a cyclic
group vanishes.  The next construction specifies the larger group
on which an alternating form does exist.

\section{Alternating doubles of the arithmetic presentation}

Return to arbitrary $r,s>0$ and $m,n,b_0\in\mathbb Z$, without
local coprimality assumptions.  Regard the actual relation matrix
\[
 A=\begin{pmatrix}r&s\\-m&n-sb_0\end{pmatrix}:\mathbb Z^2\to\mathbb Z^2
\]
as an integral map.  On $L_A=\mathbb Z^2\oplus\mathbb Z^2$, define
\begin{equation}\label{eq:doubled-form}
 \Omega_A((u,v),(u',v'))=u^TAv'-(u')^TAv,\qquad
 [\Omega_A]=\begin{pmatrix}0&A\\-A^T&0\end{pmatrix}.
\end{equation}
This construction uses the full presentation matrix and both copies of
$\mathbb Z^2$ as specified.  It is an alternating lattice associated
with the presentation, independently of a geometric fibre lattice.

\begin{theorem}[Arithmetic discriminant and its pairing]\label{thm:double}
Put $d=\gcd(r,s,m,n)$ and $\Delta=\det A$.
If $\Delta\ne0$, set $N=|\Delta|/d$.  Then
\[
 (L_A,\Omega_A)\cong (\mathbb Z^4,dJ\oplus NJ),\qquad
 \mathcal D_A\cong(\mathbb Z/d\mathbb Z)^2
                    \oplus(\mathbb Z/N\mathbb Z)^2.
\]
Here $\mathcal D_A=L_A^\vee/\Omega_A^\flat(L_A)$ and
$\Omega_A^\flat(z)=\Omega_A(z,-)$.  In the coordinates
\[
 \mathcal D_A=\operatorname{coker}A\oplus\operatorname{coker}A^T,
\]
its perfect alternating pairing is
\begin{equation}\label{eq:arithmetic-pairing}
 b_A((\bar x,\bar y),(\bar x',\bar y'))
     =x^TA^{-T}y'-(x')^TA^{-T}y\pmod{\mathbb Z}.
\end{equation}
Its exponent is $N$ and its order is $\Delta^2$.
If $\Delta=0$, the radical of $\Omega_A$ is primitive of rank two,
and its quotient lattice has form $dJ$ and discriminant group
$(\mathbb Z/d\mathbb Z)^2$.
\end{theorem}

\begin{proof}
Choose unimodular matrices $P,Q$ with $PAQ=\operatorname{diag}(d,N)$
when $\Delta\ne0$.  The change of basis
$\operatorname{diag}(P^T,Q)$ in \eqref{eq:doubled-form}, followed by
interleaving the two factors, gives $dJ\oplus NJ$.
The integer $d$ divides $N$ because $d^2$ divides $\det A$.
In the original coordinates the matrix of $\Omega_A^\flat$ is
$-[\Omega_A]$.  Its inverse over $\mathbb Q$ is
\[
        \begin{pmatrix}0&A^{-T}\\-A^{-1}&0\end{pmatrix},
\]
which gives \eqref{eq:arithmetic-pairing} by (2.1).
More explicitly, $x^TA^{-T}y$ is integral on either relation lattice:
replace $x$ by $x+Az$ or $y$ by $y+A^Tw$.
If it is integral for every integral $y$, then $A^{-1}x$ is integral,
so $x\in A\mathbb Z^2$.  The same argument applies to $y$.
This proves perfectness directly and determines the stated orders.

If $\Delta=0$, the matrix $A$ has rank one since $r,s>0$.
Its Smith form is $\operatorname{diag}(d,0)$, so the same change of
basis gives $dJ\oplus0_2$.  The last summand is the radical and is a
direct summand of $L_A$.  Taking its quotient proves the last assertion.
The unreduced cokernel also has a free summand $\mathbb Z^2$.
It is not the finite discriminant group of a nondegenerate lattice.
\end{proof}

Let $K=\operatorname{coker}A$ when $\Delta\ne0$.
The character
\[
 \chi_{\bar y}(\bar x)=\exp(2\pi i x^TA^{-T}y)
\]
identifies $\operatorname{coker}A^T$ with $\widehat K$.
Thus \eqref{eq:arithmetic-pairing} is the standard alternating pairing
on $K\oplus\widehat K$.
Define
\[
 H_A=\mu_N\times K\times\widehat K,\qquad
 (z,u,\eta)(z',u',\eta')=(zz'\eta'(u),u+u',\eta\eta').
\]
On $\mathbb C[K]$, the formula
\[
                 [\rho(z,u,\eta)h](v)=z\eta(v)h(v+u)
\]
gives a representation of dimension $|\Delta|$.
Its central character is tautological, its commutator is
$\exp(2\pi i b_A)$, and its center is $\mu_N$.
Indeed, perfectness forces any central element to have $u=0$ and
$\eta=1$.  Character averaging projects onto each delta-function,
and translations act transitively on those functions.
This proves irreducibility exactly as in Section~2, for every finite
abelian $K$, including groups of even exponent.
For $\Delta=0$, the nondegenerate quotient instead gives
$K=\mathbb Z/d\mathbb Z$ and a representation of dimension $d$.

For the non-effective data $(2,4;2,2;0)$, one has
$d=2$, $\Delta=12$, and alternating elementary divisors $(2,6)$.
Thus $\mathcal D_A\cong(\mathbb Z/2)^2\oplus(\mathbb Z/6)^2$,
and the representation has dimension twelve and central order six.
For $(2,4;2,-4;0)$, the radical has rank two and the quotient has
form $2J$.  The resulting finite group is $(\mathbb Z/2)^2$.
These are arithmetic constructions.  Nonprimitive filling slopes do
not bound meridional disks of a solid torus, so Theorem~\ref{thm:linking}
does not apply to these unreduced data.

\begin{proposition}[Separation from the level-six fibre]\label{prop:separation}
For effective two-fibre data with $(r,s)=(3,4)$ and $b_0=0$, the
arithmetic double has elementary divisors $(1,|p|)$, and
$\gcd(p,12)=1$.  In particular, it never has elementary divisors
$(1,6)$.  When $|p|=1$, its discriminant group is zero, while the
fibre lattice of Theorem~1.1 has discriminant group $(\mathbb Z/6)^2$.
\end{proposition}

\begin{proof}
Here $d=1$ and $\Delta=p=4m+3n$.
Effectiveness makes $m$ a unit modulo three and $n$ a unit modulo four.
Reduction of $p$ modulo three and four proves $\gcd(p,12)=1$.
Theorem~\ref{thm:double} now gives the result.
\end{proof}

The abstract data $(3,4;3,-2;0)$ do have $\Delta=6$ and $d=1$.
Their doubled lattice is therefore isometric to $J\oplus6J$, but
both local pairs are non-effective.  This is an equality of abstract
lattice types.  It supplies no map to Engel's fibre or to its Leray sheaf.
